% =============================================================================
% sections/design/drone/avionics_flight_control.tex  -  Avionics and Flight Control
% Teams: drone (best guess - unconfirmed) / Status: EXAMPLE (scope approximation)
%
% EXAMPLE CONTENT for scope/layout approximation - not final report text.
% Adapted from the Preliminary Report's "Core Computing and Flight Control"
% / "System Redundancy" subsections, trimmed to minimal prose. Marked with
% \exampleflag on its heading in the assembler.
%
% No separate figure - refers back to fig:drone_electronics_architecture
% (Drone Overview) rather than duplicating it.
% =============================================================================

The flight controller handles real-time control (attitude stabilization,
sensor fusion, failsafes) and directly processes \qty{2.4}{\giga\hertz}
receiver input in manual mode, while the companion computer runs heavy
autonomy tasks (computer vision, trajectory generation, state machine) and,
in autonomous mode, sends motion commands that the flight controller
safety-checks before motor actuation
(\autoref{fig:drone_electronics_architecture}).

To maximize reliability, the architecture incorporates the following
redundancies: triple-redundant IMUs in the flight controller, using a
voting algorithm to immediately isolate faulty data; a
\qty{433}{\mega\hertz} telemetry link backing up the primary
\qty{2.4}{\giga\hertz} RC link, with automatic failover verified via
mid-flight signal-loss testing; and an onboard video switcher that toggles
between the forward and downward cameras, ensuring continuous visual
feedback if one perspective is compromised.

